Dado $A$, un lenguaje sobre $\Sigma$, se define $A^R$ de la siguiente forma:
\begin{center}
    $A^R = \{u^R : u \in A\}$,
\end{center}
donde:\\
$u^R$ es la cadena inversa de $u$.\\
$A^R$ se denomina la reflexión o el inverso de $A$.

\textbf{Propiedades:} Sean $A$ y $B$ lenguajes sobre $\Sigma$ (es decir, $A,B \subseteq \Sigma^*$). Demostrar la propiedad 66 y 70.

\begin{enumerate}
    \setcounter{enumi}{65}
    \item $(AB)^R = B^RA^R$.
        \subitem Sea $A = \{a\}$ y $B = \{b\}$.
        
        Tenemos que la concatenación de $A$ y $B$ tendría la forma:
        \begin{center}
            $AB = \{ab\}$
        \end{center}
        
        Cuya forma inversa sería:
        \begin{center}
            $(AB)^R = \{ba\}$
        \end{center}
        
        Si ahora tomamos el lado derecho de la igualdad y lo desarrollamos, tenemos que:
        \begin{center}
            $B^R = \{b\}$\\
            $A^R = \{a\}$
        \end{center}
        
        Cuya concatenación es:
        \begin{center}
            $B^RA^R = \{ba\}$
        \end{center}
        
        Que es igual a la concatenación del lado izquierdo de la igualdad.
        
        Por lo tanto:
        \begin{center}
            $\boxed{(AB)^R = B^R A^R}$
        \end{center}
    %
    \setcounter{enumi}{69}
    \newpage
    \item $(A^*)^R = (A^R)^*$
        \subitem Sea $A = {a}$.
        
        Tenemos:
        \begin{center}
            $A^* = \{\lambda,a,aa,aaa,aaaa,\ldots\}$
        \end{center}
        
        Cuya forma inversa es:
        \begin{center}
            $(A^*)^R = \{a...a,\ldots,aaaa,aaa,aa,a,\lambda\}$
        \end{center}
        
        Si validamos del lado derecho de la igualdad tenemos:
        \begin{center}
            $A^R = \{a\}$
        \end{center}
        
        Cuya clausura de Kleene es:
        \begin{center}
            $(A^R)^* = \{a...a,\ldots,aaaa,aaa,aa,a,\lambda\}$
        \end{center}
        
        Por lo tanto, podemos ver que la igualdad se conserva:
        \begin{center}
            $\boxed{(A^*)^R = (A^R)^*}$
        \end{center}
    %
\end{enumerate}